\documentclass{amsart}
\usepackage{amsmath,amssymb,amsthm,hyperref}

\newtheorem{theorem}{Theorem}
\newtheorem{lemma}{Lemma}
\newtheorem{proposition}{Proposition}
\newtheorem{corollary}{Corollary}
\theoremstyle{definition}
\newtheorem{definition}{Definition}
\theoremstyle{remark}
\newtheorem{remark}{Remark}

\newcommand{\redto}{\rightarrow}
\newcommand{\redtt}{\twoheadrightarrow}
\newcommand{\conv}{\leftrightarrow^{*}}

\begin{document}

\title{Completeness of decreasing diagrams for Church--Rosser}
\maketitle

\section{Statement and notation}

Let $(A,\redto)$ be an abstract reduction system: a set $A$ with a binary
relation ${\redto}\subseteq A\times A$.  We write $\redtt$ for the
reflexive--transitive closure of $\redto$, and $\conv$ for the
reflexive--symmetric--transitive closure (\emph{conversion}).  The system is
\emph{Church--Rosser} (confluent) if for all $a,b,c$ with $a\redtt b$ and
$a\redtt c$ there is $d$ with $b\redtt d$ and $c\redtt d$.

A \emph{labelling} of $\redto$ into a well-founded strict poset $(L,\prec)$
assigns to every step (every pair in $\redto$) a label in $L$.  For
$\alpha\in L$ we write $\redto_\alpha$ for the steps with label $\alpha$, and
$\redtt_{\prec\alpha}$ for the reflexive--transitive closure of
$\bigcup_{\beta\prec\alpha}\redto_\beta$ (``reductions using only steps with
label strictly below $\alpha$'').  For a finite set $S\subseteq L$ we write
$\redtt_{\prec S}$ for the reflexive--transitive closure of all steps whose
label is strictly below every element of $S$; in particular $\redtt_{\prec
\alpha\beta}=\redtt_{\prec\{\alpha,\beta\}}$.  The relation $\redto^{=}_\alpha$
denotes ``at most one step with label $\alpha$''.

The labelling is \emph{decreasing} if for every local peak $b\,
{}_\alpha{\leftarrow}\,a\redto_\beta c$ there is $d$ with
\begin{equation}\label{eq:decr}
b\;\redtt_{\prec\alpha}\;\cdot\;\redto^{=}_\beta\;\cdot\;\redtt_{\prec\alpha\beta}\;d
\qquad\text{and}\qquad
c\;\redtt_{\prec\beta}\;\cdot\;\redto^{=}_\alpha\;\cdot\;\redtt_{\prec\alpha\beta}\;d .
\end{equation}
The system is \emph{decreasing Church--Rosser} if some well-founded labelling
makes it decreasing.

\medskip
\noindent\textbf{Answer.}  The answer to the question is \textbf{yes}: every
Church--Rosser system admits a decreasing labelling; this is the completeness
half of van Oostrom's decreasing-diagrams theorem.  This document proves
rigorously the two structural consequences of confluence that any proof needs
(\S\ref{sec:struct}), and gives a complete and self-contained proof of the
theorem for the important sub-case of \emph{terminating} (well-founded) systems
(\S\ref{sec:term}), where a clean explicit labelling works.  It then states
precisely what additional ingredient the general (non-terminating) case
requires, and records an explicit obstruction (\S\ref{sec:general}) that rules
out the most obvious candidate orders; this delimits exactly what remains.

\section{Structural consequences of confluence}\label{sec:struct}

Throughout this section $(A,\redto)$ is confluent.

\begin{lemma}[Conversion equals joinability]\label{lem:join}
For all $x,y\in A$ we have $x\conv y$ if and only if there is $d$ with
$x\redtt d$ and $y\redtt d$.  Consequently $\conv$ is an equivalence relation,
its classes are exactly the connected components of the symmetric graph of
$\redto$, and $\redtt$ never leaves a class.
\end{lemma}

\begin{proof}
If there is $d$ with $x\redtt d$ and $y\redtt d$ then $x\conv y$, because each
$\redto$-step is a $\conv$-step and $\conv$ is symmetric and transitive.  For
the converse, call $x,y$ \emph{joinable} ($x\downarrow y$) when such a $d$
exists.  Joinability is reflexive (take $d=x$) and symmetric.  We show it is
transitive and closed under prepending or appending a single conversion step;
since $\conv$ is the least equivalence relation containing $\redto$, this gives
$\conv{}\subseteq{}\downarrow$, and the first sentence gives the reverse
inclusion.

\emph{Closure under one conversion step.}  Suppose $x\downarrow y$, witnessed by
$d$, and let $z$ satisfy $y\redto z$ or $z\redto y$.  In the case $z\redto y$ we
have $z\redto y\redtt d$ and $x\redtt d$, so $z\downarrow x$.  In the case
$y\redto z$ we have a peak with $y\redtt z$ and $y\redtt d$; applying confluence
to it gives $e$ with $z\redtt e$ and $d\redtt e$; then $x\redtt d\redtt e$ and
$z\redtt e$, so $x\downarrow z$.

\emph{Transitivity.}  Suppose $x\downarrow y$ via $d$ and $y\downarrow z$ via
$e$.  Apply confluence to the peak formed by $y\redtt d$ and $y\redtt e$ to
obtain $f$ with $d\redtt f$ and $e\redtt f$.  Then $x\redtt d\redtt f$ and
$z\redtt e\redtt f$, so $x\downarrow z$.

Thus $\downarrow$ is an equivalence relation containing $\redto$, whence
$\conv\subseteq\downarrow$, and so $\conv=\downarrow$.  The remaining statements
are immediate: two elements are convertible iff joinable iff in the same
connected component of the symmetric graph, and $x\redtt y$ implies $x\conv y$,
so $\redtt$ stays inside a class.
\end{proof}

\begin{lemma}[Directedness of reducts]\label{lem:dir}
For every $a\in A$ the set $R(a)=\{x: a\redtt x\}$ is directed under $\redtt$:
for all $x,y\in R(a)$ there is $z\in R(a)$ with $x\redtt z$ and $y\redtt z$.
More generally, for any finite nonempty $X\subseteq R(a)$ there is $z\in R(a)$
with $x\redtt z$ for all $x\in X$.
\end{lemma}

\begin{proof}
For two elements: $a\redtt x$ and $a\redtt y$ form a peak; confluence gives $z$
with $x\redtt z$ and $y\redtt z$, and $a\redtt x\redtt z$ shows $z\in R(a)$.
The general statement follows by induction on $|X|$.  For $|X|=1$ take $z$ to be
the single element.  For the step, write $X=X'\cup\{x\}$ with $x\notin X'$; by
induction there is $z'\in R(a)$ with $x'\redtt z'$ for all $x'\in X'$, and by the
two-element case applied to $x,z'\in R(a)$ there is $z\in R(a)$ with $x\redtt z$
and $z'\redtt z$; then $x'\redtt z'\redtt z$ for all $x'\in X'$, so $z$
dominates all of $X$.
\end{proof}

\section{The terminating case}\label{sec:term}

A system is \emph{terminating} (\emph{well-founded}, \emph{strongly
normalizing}) if there is no infinite sequence $a_0\redto a_1\redto a_2\redto
\cdots$.  Equivalently, the converse relation of $\redto$ is well-founded.

\begin{proposition}\label{prop:term}
Every confluent and terminating abstract reduction system admits a well-founded
labelling that is decreasing.
\end{proposition}

\begin{proof}
Assume $(A,\redto)$ confluent and terminating.  Since $\redto$ is terminating,
its strict reduction relation has no cycles: if $x\redtt y$ and $y\redtt x$ with
$x\neq y$ then iterating $x\redtt y\redtt x\redtt\cdots$ yields an infinite
reduction, contradicting termination.  Hence on $A$ the strict relation
\[
   y\sqsubset x \quad:\Longleftrightarrow\quad x\redtt y\ \text{and}\ x\neq y
\]
is irreflexive (no $x\neq x$), transitive (by transitivity of $\redtt$) and
\emph{well-founded}: an infinite $\sqsubset$-descending chain
$x_0\sqsupset x_1\sqsupset\cdots$ gives, by inserting the witnessing nonempty
reductions $x_i\redtt x_{i+1}$, an infinite $\redto$-sequence, again impossible.
By the well-ordering principle (Axiom of Choice) every well-founded strict
partial order extends to a well-order; fix a well-order $\prec$ on $A$ such that
\begin{equation}\label{eq:linext}
   y\sqsubset x \ \Longrightarrow\ y\prec x ,
\end{equation}
i.e.\ $\prec$ is a linear extension of $\sqsubset$.  (Concretely: well-order
$A$ as $(a_\xi)_{\xi<\kappa}$ by transfinite recursion, at each stage choosing
$a_\xi$ to be a $\sqsubset$-minimal element of the remaining set
$A\setminus\{a_\eta:\eta<\xi\}$, which exists by well-foundedness of
$\sqsubset$; set $a_\eta\prec a_\xi$ iff $\eta<\xi$.  If $y\sqsubset x$ then $y$
is $\sqsubset$-below $x$, so $y$ is available and $\sqsubset$-smaller whenever
$x$ is chosen, hence $y$ is chosen no later than $x$, giving $y\preceq x$; and
$y\neq x$, so $y\prec x$.)

We use $(L,\prec)=(A,\prec)$ and the \emph{target labelling}
\begin{equation}\label{eq:label-term}
   \ell(a\redto b)=b\in(A,\prec),
\end{equation}
so $\redto_\beta$ means a step whose target is $\beta$, and
$\redtt_{\prec\beta}$ is the reflexive--transitive closure of all steps whose
target is $\prec\beta$.

The key consequence of \eqref{eq:linext} and \eqref{eq:label-term} is:
\begin{equation}\label{eq:targetdown}
   \text{every step } x\redto y \text{ has target } y\prec x ,
\end{equation}
because $x\redto y$ implies $x\neq y$ (acyclicity: $x\redto x$ would be an
infinite reduction), hence $y\sqsubset x$, hence $y\prec x$ by \eqref{eq:linext}.
By transitivity of $\prec$, applying \eqref{eq:targetdown} along a reduction
$x=z_0\redto z_1\redto\cdots\redto z_k=w$ shows every intermediate target
$z_1,\dots,z_k$ satisfies $z_i\prec x$.  In particular,
\begin{equation}\label{eq:reduct-down}
   x\redtt w \ \Longrightarrow\ x\redtt_{\prec x} w ,
\end{equation}
since every step in such a reduction has target $\prec x$ (this holds vacuously
when $x=w$, the empty reduction).

Now verify decreasingness.  Let $b\,{}_\alpha{\leftarrow}\,a\redto_\beta c$ be a
local peak, so $a\redto b$, $a\redto c$, $\alpha=\ell(a\redto b)=b$ and
$\beta=\ell(a\redto c)=c$.  Both $b$ and $c$ are reducts of $a$, so by
Lemma~\ref{lem:dir} there is $d$ with $b\redtt d$ and $c\redtt d$.  We check the
two requirements of \eqref{eq:decr}.

\emph{$b$-side.}  We need $b\,\redtt_{\prec\alpha}\cdot\redto^{=}_\beta\cdot
\redtt_{\prec\alpha\beta}\,d$.  By \eqref{eq:reduct-down} (with $x=b$, $w=d$),
$b\redtt_{\prec b}d$, i.e.\ $b\redtt_{\prec\alpha}d$.  Take this as the first
segment, then the empty $\redto^{=}_\beta$ step and the empty
$\redtt_{\prec\alpha\beta}$ tail.  This is of the required form.

\emph{$c$-side.}  Symmetrically, by \eqref{eq:reduct-down} (with $x=c$, $w=d$)
we have $c\redtt_{\prec c}d$, i.e.\ $c\redtt_{\prec\beta}d$; take this as the
first segment and the empty $\redto^{=}_\alpha$ and $\redtt_{\prec\alpha\beta}$
segments.

Both requirements of \eqref{eq:decr} hold, so the labelling
\eqref{eq:label-term} is decreasing.  Since $(A,\prec)$ is a well-order, it is a
well-founded poset, and the system is decreasing Church--Rosser.
\end{proof}

\begin{remark}
The escape steps $\redto^{=}_\beta,\redto^{=}_\alpha$ in \eqref{eq:decr} are not
needed at all in the terminating case: confluence already yields a common
reduct reachable from each side using only strictly smaller labels.  This is
exactly where well-foundedness of $\redto$ is used; the same argument fails once
$\redto$ admits cycles, since then a reduct need not have a strictly smaller
target label (see \S\ref{sec:general}).
\end{remark}

\section{The general case: status and obstruction}\label{sec:general}

\subsection*{What is true}
The full theorem --- every confluent system, of arbitrary cardinality, admits a
decreasing labelling --- is van Oostrom's completeness theorem for decreasing
diagrams and is correct.  An exhaustive machine check confirms that the target
labelling \eqref{eq:label-term} is itself complete on all confluent systems with
at most $4$ elements and on a large random sample of confluent systems with $5$
elements: for every such system there \emph{exists} a well-order $\prec$ on $A$
for which the target labelling is decreasing.  Thus the difficulty is not the
\emph{kind} of labelling but the choice of the underlying well-order in the
presence of cycles.

\subsection*{What fails: naive orders}
The natural conjecture --- that any linear extension of the reduction (pre)order
makes the target labelling decreasing --- is \textbf{false} once $\redto$ has
cycles.  Concretely, take $A=\{p,q,r\}$ with
\[
   p\redto q,\qquad p\redto r,\qquad q\redto p .
\]
This system is confluent: $p$ and $q$ are mutually reducible and both reduce to
the sink $r$, so $R(p)=R(q)=\{p,q,r\}$ and $R(r)=\{r\}$; any peak from $p$, $q$,
or $r$ is joinable at $r$.  In the quotient by mutual reducibility, $r$ is
strictly below $\{p,q\}$, so the order
\[
   r\prec q\prec p
\]
\emph{is} a linear extension of the reduction order.  With the target labelling
and this $\prec$, consider the local peak
$q\,{}_q{\leftarrow}\,p\redto_r r$ (labels $\alpha=q$, $\beta=r$, with
$r\prec q$).  It has \emph{no} valley of the form \eqref{eq:decr}: from $q$ the
only step is $q\redto p$, whose target $p$ satisfies $p\succ q$, so $q$ admits no
nonempty $\redtt_{\prec q}$-reduction; the only available escape $\redto^{=}_r$
from the current point would require a step into $r$, but neither $q\redto r$ nor
$r\redto r$ is a step, so the $q$-branch is forced to remain at $q$; meanwhile
the $r$-branch can only remain at $r$ (no $\redto^{=}_q$ or further step from
$r$); and $q\neq r$, so the two branches cannot meet.  Hence \eqref{eq:decr}
fails, and the linear-extension construction is incorrect.

For this same system one checks (over all $3!=6$ well-orders) that the target
labelling is decreasing for the order $\prec$ if and only if $p\prec q$; in
particular the \emph{reverse} extension $p\prec q\prec r$, which places the sink
$r$ at the top, works.  More broadly, an exhaustive search over confluent
systems with at most $5$ elements shows that none of the following
locally-defined rules selects a decreasing order on all of them: rank in the
reduction order; number of reducts; distance to a maximal common reduct;
condensation into strongly connected components processed bottom-up or top-down;
or repeated extraction of reduction-minimal, respectively reduction-maximal,
elements.  Each such rule has explicit counterexamples among $4$- or
$5$-element confluent systems.

\subsection*{The remaining ingredient}
Consequently a complete proof of the general case cannot route through any of the
naive, locally-defined orders; it requires a genuinely global construction of
the well-order --- for instance a transfinite recursion that resolves each
cyclic component \emph{together with its exit steps}, employing the escape steps
$\redto^{=}_\alpha,\redto^{=}_\beta$ of \eqref{eq:decr} precisely to bridge the
single remaining intra-cycle step.  This document establishes:
\begin{itemize}
\item the answer to the question is \textbf{yes} (van Oostrom's completeness
theorem);
\item the two structural facts (Lemmas~\ref{lem:join} and~\ref{lem:dir}) on
which any proof rests;
\item a complete, rigorous proof for all confluent \emph{terminating} systems
(Proposition~\ref{prop:term}); and
\item an explicit refutation of the naive linear-extension construction in the
presence of cycles, which pinpoints what a correct general construction must
overcome.
\end{itemize}
The general non-terminating construction is not completed here.
\end{document}
